% Generated draft: requires mathematical and source review.
\documentclass[11pt]{article}
\usepackage[margin=25mm]{geometry}
\usepackage{fontspec}
\setmainfont{Latin Modern Roman}
\usepackage{amsmath,amssymb,mathtools,unicode-math}
\setmathfont{Latin Modern Math}
\usepackage{enumitem,needspace,xurl,hyperref}
\setlength{\parindent}{0pt}
\setlength{\parskip}{5pt}
\setlength{\emergencystretch}{3em}
\newcommand{\N}{\mathbb{N}}
\newcommand{\Z}{\mathbb{Z}}
\newcommand{\Q}{\mathbb{Q}}
\newcommand{\R}{\mathbb{R}}
\newcommand{\C}{\mathbb{C}}
\newcommand{\F}{\mathbb{F}}
\newcommand{\A}{\mathbb{A}}
\newcommand{\PP}{\mathbb{P}}
\newcommand{\E}{\mathbb{E}}
\newcommand{\eps}{\varepsilon}
\newcommand{\dd}{\,\mathrm{d}}
\begin{document}
\section*{UnsolvedMath ID: 11100061}
\textbf{Source problem code:} catalogue entry
\subsection{Definitions and mathematical statement}
Hovey asks to determine \(v_1\)-exponents of spheres, while explicitly saying that he is not sure of the precise statement or an expected conjectural formula. The motivating comparison is the reported Cohen--Moore--Neisendorfer result: for an odd prime \(p\) and \(n\geq1\), the \(p\)-primary torsion in \(\pi_*(S^{2n+1})\) is annihilated by \(p^n\), with \(p^{n-1}\) insufficient. This is the standard exponent reading of the source's plain-text ``not p\^{}n-1,'' not an assertion about the nontorsion fundamental sphere class.

The source suggests maps from a Moore space so an Adams self-map acts. For example, after making the necessary suspension and prime choices, let \(V\) be a based Moore space and \(v:\Sigma^dV\to V\) an Adams \(v_1\)-self-map. Precomposition gives maps
\[v^*:[\Sigma^jV,X]\to[\Sigma^{j+d}V,X].\]
The Moore exponent, unstable realization of the self-map and group/degree range must be fixed; a stable self-map alone does not automatically define every proposed unstable operation.

There are at least two materially different exponent questions. One can consider \(v\)-torsion classes (killed by some iterate of \(v^*\)) and ask for a uniform power \((v^*)^N\) killing all of them. Alternatively one can invert \(v\) in the directed system and ask for the \(p\)-power exponent of the resulting periodic groups. The first is a \(v\)-torsion exponent and the second a \(p\)-exponent after periodic localization; they are not interchangeable. The summary adopts the latter language, but Hovey's uncertain text does not decide between them, nor specify which Moore spaces/coefficients should be used. Accordingly no one of these is silently chosen as the problem, and no numeric formula is inserted.

The broad sphere range, possible prime-two version and expected bound remain unspecified; only the motivating ordinary-torsion comparison explicitly concerns odd primes. The final source request is for clarification and an expected conjecture, not a declared value.
\subsection{Short English statement}
Clarify and determine the intended \(v_1\)-exponent of spheres using Moore-space Adams self-maps. The source itself leaves open whether this means a \(v_1\)-torsion bound or a \(p\)-exponent of periodic groups.
\subsection{Sources}
\begin{itemize}
\item Local UnsolvedMath registry, permanent ID \(11100061\), exact source code \texttt{catalogue entry}.
\item Mark Hovey, \emph{Hovey's Algebraic Topology Problem List}, Unstable homotopy theory, item 2, surviving mirror: \url{https://www-users.cse.umn.edu/~tlawson/hovey/}. Registry access 2026-07-29; html-ordered-list extraction.
\item Cohen--Moore--Neisendorfer are credited for the ordinary odd-primary exponent comparison. The summary mentions Bousfield, Davis--Mahowald, Toda and a truncated Davis lead without precise citations; its self-corrected names and placeholder exponent expression are not usable theorem statements.
\end{itemize}
\end{document}
